\section{Composição de funções contínuas e limites}

Nessa seção, a menos que explicitamente indicado, consideraremos que $(M, d_M)$, $(N, d_N)$ e $(P, d_P)$ são espaços métricos arbitrários.

Estudaremos como se comportam limites a luz da composição de funções.

O teorema abaixo diz que, sob hipóteses adequadas, a continuidade de funções permite que o limite ``entre'' em uma composição.

\begin{proposition}
    Sejam
    \begin{itemize}
        \item $A\subseteq M$ um subespaço,
        \item $B\subseteq N$ um subespaço,
        \item $f: A \to B$ uma função, e
        \item $g: B \to P$ uma função.
        \item $a \in M$ um ponto de acumulação de $A$.
        \item $L \in B$.
    \end{itemize}

    Suponha que $L$ é limite de $f$ em $a$.
    Então $g(L)$ é limite de $g\circ f$ em $a$.

    Em outras palavras, se $\lim_{x\to a} f(x)$ existe e $g$ é uma função contínua em $\lim_{x\to a} f(x)$, então:

    \[
        \lim_{x\to a} g(f(x)) = g\left(\lim_{x\to a} f(x)\right).
    \]
\end{proposition}

\begin{proof}
    Suponha que $L$ é limite de $f$ em $a$.
    Devemos mostrar que $g(L)$ é limite de $g\circ f$ em $a$.

    Devemos ver que, fixado $\epsilon>0$, existe $\delta>0$ tal que, para todo $x \in A\setminus\{a\}$, 
    \[
        \text{se } d_M(x, a) < \delta,\qquad \text{então} \qquad d_P(g(f(x)), g(L)) < \epsilon.
    \]

    Fixe $\epsilon>0$.
    Como $g$ é contínua em $L$, existe $\nu>0$ tal que, para todo $y \in B$, se $d_N(y, L) < \nu$, então $d_P(g(y), g(L)) < \epsilon$.

    Como $L$ é limite de $f$ em $a$, existe $\delta>0$ tal que, para todo $x \in A\setminus\{a\}$, se $d_M(x, a) < \delta$, então $d_N(f(x), L) < \nu$.
    
    O número $\delta>0$ acima é como buscado: suponha que $x \in A\setminus\{a\}$ é tal que $d_M(x, a) < \delta$.
    Então $d_N(f(x), L) < \nu$.
    Pondo $y = f(x)$, temos que $d_N(y, L) < \nu$, e, portanto, $d_P(g(y), g(L)) < \epsilon$.
\end{proof}

\begin{corollary}
    Sejam
    \begin{itemize}
        \item $f: M \to N$ uma função, e
        \item $g: N \to P$ uma função.
        \item $a \in M$ um ponto de acumulação de $M$.
    \end{itemize}

    Se $f$ é contínua em $a$ e $g$ é contínua em $f(a)$, então $g\circ f$ é contínua em $a$.

    Consequentemente, a composição de funções contínuas é contínua.
\end{corollary}

\begin{proof}
    Se $a \in M$ é um ponto isolado, então trivialmente $g\circ f$ é contínua em $a$.

    Se $a \in M$ é um ponto de acumulação de $M$, então
    \[
        \lim_{x\to a} f(x) = f(a),
    \]

    Assim, da proposição anterior, temos

    \[
        \lim_{x\to a} g(f(x)) = g\left(\lim_{x\to a} f(x)\right) = g(f(a)).
    \]

    Portanto, $g\circ f$ é contínua em $a$.
\end{proof}

O teorema abaixo nos mostra que, sob certas hipóteses, podemos ``compor limites''.
Alguns textos encaram o teorema abaixo como um teorema de substituição de variáveis em limites.

\begin{proposition}
    Sejam
    \begin{itemize}
        \item $A\subseteq M$ um subespaço,
        \item $B\subseteq N$ um subespaço,
        \item $f: A \to B$ uma função, e
        \item $g: B \to P$ uma função.
        \item $a \in M$ um ponto de acumulação de $A$.
        \item $b \in N$ um ponto de acumulação de $B$.
        \item $L \in P$.
    \end{itemize}


    Suponha que:
    \begin{itemize}
        \item $L$ é limite de $g$ em $b$, e
        \item $b$ é limite de $f$ em $a$.
        \item $f(x)\neq b$ para todo $x \in A\setminus\{a\}$.
    \end{itemize}

    Então $L$ é limite de $g\circ f$ em $a$.
    Em outras palavras, nessas hipóteses,
    \[
        \lim_{x\to a} g(f(x)) = \lim_{y\to b} g(y).
    \]
\end{proposition}

\begin{proof}
    Devemos mostrar que $L$ é limite de $g\circ f$ em $a$.

    Devemos ver que, fixado $\epsilon>0$, existe $\delta>0$ tal que, para todo $x \in A\setminus\{a\}$, 
    \[
        \text{se } d_M(x, a) < \delta,\qquad \text{então} \qquad d_P(g(f(x)), L) < \epsilon.
    \]

    Fixe $\epsilon>0$.
    Como $L$ é limite de $g$ em $b$, existe $\nu>0$ tal que, para todo $y \in B$, se $d_N(y, L) < \nu$, então $d_P(g(y), g(L)) < \epsilon$.

    Como $b$ é limite de $f$ em $a$, existe $\delta>0$ tal que, para todo $x \in A\setminus\{a\}$, se $d_M(x, a) < \delta$, então $d_N(f(x), b) < \nu$.

    O número $\delta>0$ acima é como buscado: suponha que $x \in A\setminus\{a\}$ é tal que $d_M(x, a) < \delta$.
    Então $d_N(f(x), b) < \nu$ e $f(x)\in B\setminus\{b\}$.
    Segue que $d_P(g(f(x)), g(L)) < \epsilon$.
\end{proof}